% GENERATED FILE. Edit canonical lessons, then run npm run generate:books.
% canonical-source-hash: fnv1a64:2b1498d34015ba13
% canonical-lessons: JA-C321-zettai, JA-C321-kanarazu, JA-C321-yahari, JA-C321-mushiro, JA-C321-unazuku

\chapter{How Sure Are You?}
\label{ch:ja-how-sure-are-you}
\hlchaptermodality{321}

\begin{chapteropening}
\textbf{By the end of this chapter:} \emph{I can say how sure I am, and why.}

Complete the last lesson of chapter 321: I can say how sure I am, and why.
\end{chapteropening}

\section[\texorpdfstring{\ja{ぜったい} (zettai)}{ぜったい (zettai)}]{\ja{ぜったい} (zettai) — absolutely; definitely}

\label{lesson:JA-C321-zettai}



\begin{quote}
Before the new one: say the Japanese for to comprehend, then the Japanese for especially.
\end{quote}

\subsection*{You'll want to know: \ja{ぜったい}}
\textbf{\ja{ぜったい}} — \emph{zettai} — \textquotedblleft{}absolutely; definitely\textquotedblright{}.

\begin{quote}
\textbf{\ja{あぶないから、ぜったい} \ja{いかないで} \ja{ください。}} — \emph{abunai kara, zettai ikanaide kudasai} — \textquotedblleft{}It's dangerous, so do not go there under any circumstances.\textquotedblright{}
\end{quote}

\begin{cousinweb}
\textbf{\ja{ぜったい}} (\emph{zettai}) is Sino-Japanese (kango): two Chinese characters read with their on'yomi, readings borrowed from Chinese. They are \emph{zetsu}, \textquotedblleft{}cut off\textquotedblright{} (Middle Chinese \emph{dzjwet}) and \emph{tai}, \textquotedblleft{}opposite\textquotedblright{} (\emph{twojH}) (KANJIDIC2; Middle Chinese in Baxter's transcription). The two join as \emph{zet-tai}.
\end{cousinweb}

\begin{grammarlens}[title={What you've built}]
One more word for giving reasons and opinions.
\end{grammarlens}

\subsection*{Your turn}
\begin{itemize}
\raggedright
  \item \emph{Say it:} \emph{zettai}
  \item \emph{Say it:} \emph{zettai}, once more
  \item \emph{Say it:} \emph{abunai kara, zettai ikanaide kudasai} — It's dangerous, so do not go there under any circumstances.
  \item \emph{Recall:} say the Japanese for to comprehend, then the Japanese for especially, then say \emph{zettai} again
\end{itemize}

\begin{tcolorbox}[breakable,colback=teal!4,colframe=teal!35!black,title={Before you move on}]
What does \ja{ぜったい} mean? (\textquotedblleft{}Absolutely; definitely\textquotedblright{}.) Say it once more.
\end{tcolorbox}
\section[\texorpdfstring{\ja{かならず} (kanarazu)}{かならず (kanarazu)}]{\ja{かならず} (kanarazu) — without fail; certainly}

\label{lesson:JA-C321-kanarazu}



\begin{quote}
Before the new one: say the Japanese for especially, then the Japanese for absolutely.
\end{quote}

\subsection*{You'll want to know: \ja{かならず}}
\textbf{\ja{かならず}} — \emph{kanarazu} — \textquotedblleft{}without fail; certainly\textquotedblright{}.

\begin{quote}
\textbf{\ja{やくそくを} \ja{したから、かならず} \ja{いきます。}} — \emph{yakusoku o shita kara, kanarazu ikimasu} — \textquotedblleft{}I promised, so I will go without fail.\textquotedblright{}
\end{quote}

\begin{cousinweb}
\textbf{\ja{かならず}} (\emph{kanarazu}) is a native Japanese word (wago), inherited rather than borrowed from Chinese, and its older origin is not settled. Its kanji was attached to it later as a kun'yomi, a native reading given to a Chinese character (JMdict, KANJIDIC2).
\end{cousinweb}

\begin{grammarlens}[title={What you've built}]
One more word for giving reasons and opinions.
\end{grammarlens}

\subsection*{Your turn}
\begin{itemize}
\raggedright
  \item \emph{Say it:} \emph{kanarazu}
  \item \emph{Say it:} \emph{kanarazu}, once more
  \item \emph{Say it:} \emph{yakusoku o shita kara, kanarazu ikimasu} — I promised, so I will go without fail.
  \item \emph{Recall:} say the Japanese for especially, then the Japanese for absolutely, then say \emph{kanarazu} again
\end{itemize}

\begin{tcolorbox}[breakable,colback=teal!4,colframe=teal!35!black,title={Before you move on}]
What does \ja{かならず} mean? (\textquotedblleft{}Without fail; certainly\textquotedblright{}.) Say it once more.
\end{tcolorbox}
\section[\texorpdfstring{\ja{やはり} (yahari)}{やはり (yahari)}]{\ja{やはり} (yahari) — as expected; after all}

\label{lesson:JA-C321-yahari}



\begin{quote}
Before the new one: say the Japanese for absolutely, then the Japanese for without fail.
\end{quote}

\subsection*{You'll want to know: \ja{やはり}}
\textbf{\ja{やはり}} — \emph{yahari} — \textquotedblleft{}as expected; after all\textquotedblright{}.

\begin{quote}
\textbf{\ja{やはり} \ja{ははの} \ja{いけんが} \ja{ただしいと} \ja{おもいます。}} — \emph{yahari haha no iken ga tadashii to omoimasu} — \textquotedblleft{}After all, I think my mother's opinion is right.\textquotedblright{}
\end{quote}

\begin{cousinweb}
\textbf{\ja{やはり}} (\emph{yahari}) is a native Japanese word (wago), inherited rather than borrowed from Chinese, and its older origin is not settled. JMdict notes that it is usually written in kana (JMdict).
\end{cousinweb}

\begin{grammarlens}[title={What you've built}]
One more word for giving reasons and opinions.
\end{grammarlens}

\subsection*{Your turn}
\begin{itemize}
\raggedright
  \item \emph{Say it:} \emph{yahari}
  \item \emph{Say it:} \emph{yahari}, once more
  \item \emph{Say it:} \emph{yahari haha no iken ga tadashii to omoimasu} — After all, I think my mother's opinion is right.
  \item \emph{Recall:} say the Japanese for absolutely, then the Japanese for without fail, then say \emph{yahari} again
\end{itemize}

\begin{tcolorbox}[breakable,colback=teal!4,colframe=teal!35!black,title={Before you move on}]
What does \ja{やはり} mean? (\textquotedblleft{}As expected; after all\textquotedblright{}.) Say it once more.
\end{tcolorbox}
\section[\texorpdfstring{\ja{むしろ} (mushiro)}{むしろ (mushiro)}]{\ja{むしろ} (mushiro) — rather; if anything}

\label{lesson:JA-C321-mushiro}



\begin{quote}
Before the new one: say the Japanese for without fail, then the Japanese for as expected.
\end{quote}

\subsection*{You'll want to know: \ja{むしろ}}
\textbf{\ja{むしろ}} — \emph{mushiro} — \textquotedblleft{}rather; if anything\textquotedblright{}.

\begin{quote}
\textbf{\ja{でんしゃより} \ja{むしろ} \ja{じてんしゃが} \ja{はやいと} \ja{おもいます。}} — \emph{densha yori mushiro jitensha ga hayai to omoimasu} — \textquotedblleft{}I think a bicycle is, if anything, faster than the train.\textquotedblright{}
\end{quote}

\begin{cousinweb}
\textbf{\ja{むしろ}} (\emph{mushiro}) is a native Japanese word (wago), inherited rather than borrowed from Chinese, and its older origin is not settled. JMdict notes that it is usually written in kana (JMdict).
\end{cousinweb}

\begin{grammarlens}[title={What you've built}]
One more word for giving reasons and opinions.
\end{grammarlens}

\subsection*{Your turn}
\begin{itemize}
\raggedright
  \item \emph{Say it:} \emph{mushiro}
  \item \emph{Say it:} \emph{mushiro}, once more
  \item \emph{Say it:} \emph{densha yori mushiro jitensha ga hayai to omoimasu} — I think a bicycle is, if anything, faster than the train.
  \item \emph{Recall:} say the Japanese for without fail, then the Japanese for as expected, then say \emph{mushiro} again
\end{itemize}

\begin{tcolorbox}[breakable,colback=teal!4,colframe=teal!35!black,title={Before you move on}]
What does \ja{むしろ} mean? (\textquotedblleft{}Rather; if anything\textquotedblright{}.) Say it once more.
\end{tcolorbox}
\section[\texorpdfstring{\ja{うなずく} (unazuku)}{うなずく (unazuku)}]{\ja{うなずく} (unazuku) — to nod}

\label{lesson:JA-C321-unazuku}



\begin{quote}
Before the new one: say the Japanese for as expected, then the Japanese for rather.
\end{quote}

\subsection*{You'll want to know: \ja{うなずく}}
\textbf{\ja{うなずく}} — \emph{unazuku} — \textquotedblleft{}to nod\textquotedblright{}.

\begin{quote}
\textbf{\ja{わかったので、こどもは} \ja{うなずきました。}} — \emph{wakatta node, kodomo wa unazukimashita} — \textquotedblleft{}The child understood, so he nodded.\textquotedblright{}
\end{quote}

Polite \textbf{\ja{うなずきます}}, \emph{unazukimasu}; a godan verb. Past (polite): \textbf{\ja{こどもが} \ja{うなずきました。}}, \emph{kodomo ga unazukimashita}.

\begin{cousinweb}
\textbf{\ja{うなずく}} (\emph{unazuku}) is a native Japanese word (wago), inherited rather than borrowed from Chinese, and its older origin is not settled. Its kanji was attached to it later as a kun'yomi, a native reading given to a Chinese character (JMdict, KANJIDIC2).
\end{cousinweb}

\begin{grammarlens}[title={What you've built}]
One more word for giving reasons and opinions.
\end{grammarlens}

\subsection*{Your turn}
\begin{itemize}
\raggedright
  \item \emph{Say it:} \emph{unazuku}
  \item \emph{Say it:} \emph{unazuku}, once more
  \item \emph{Say it:} \emph{wakatta node, kodomo wa unazukimashita} — The child understood, so he nodded.
  \item \emph{Recall:} say the Japanese for as expected, then the Japanese for rather, then say \emph{unazuku} again
\end{itemize}

\begin{tcolorbox}[breakable,colback=teal!4,colframe=teal!35!black,title={Before you move on}]
What does \ja{うなずく} mean? (\textquotedblleft{}To nod\textquotedblright{}.) Say it once more.
\end{tcolorbox}
